\documentclass[11pt]{article}
\usepackage[a4paper,margin=1in]{geometry}
\usepackage{amsmath}
\usepackage{xcolor}

\definecolor{commentgray}{gray}{0.35}
\newenvironment{refcomment}%
  {\begin{quote}\color{commentgray}\itshape}%
  {\end{quote}}
\newcommand{\response}{\medskip\noindent\textbf{Response.}\ }
\newcommand{\point}[1]{\bigskip\noindent\textbf{#1}}

\setlength{\parskip}{0.4em}
\setlength{\parindent}{0pt}

\title{Response to Referee 3\\[2pt]
\large Manuscript JFM-2026-1781\\
\emph{Strain-coupled one-dimensional turbulence for rapid distortion}\\
Sharma \& Kerstein}
\date{}

\begin{document}
\maketitle

We thank the referee. Both of the fundamental objections were correct, and
addressing them has changed the theory rather than only its presentation.
The rigid-translation object is no longer treated as a linear-RDT
benchmark; the exact projected RDT, which is broadband, is now the linear
reference, and the spectral results are validated against a strained-box
DNS. Section~5 no longer claims a closed single-line functional: it states
the obstruction exactly, bounds the rapid term from line data, and measures
the axisymmetry hypothesis instead of assuming it. Established results,
assumptions and new developments are separated, the hypothesis labelled A1
is defined at first use, and the kernel equations are derived in an
appendix. Section, figure, table and equation numbers below refer to the
revised manuscript.

\point{1. The linear-RDT spectral benchmark is not justified.}
\begin{refcomment}
``the manuscript treats rigid translation of the one-dimensional spectrum
as an exact linear-RDT benchmark, although this is neither derived from the
full RDT equations nor demonstrated by a full spectral RDT calculation.
\ldots\ In Equation (2.14), the Fourier velocity amplitudes depend on the
full three-dimensional wavevector. The manuscript provides no derivation
showing that the projected spectrum must undergo a rigid translation.''
\end{refcomment}

\response
We accept this in full and have removed the incorrect identification. The
$\mathrm{d}\kappa_2/\mathrm{d}t=-A_{22}\kappa_2$ relation is now presented
for what it is---the kinematics of the ODT domain dilatation
(\S\,2.5, equation~(2.20))---and the eddy-free run that follows it is
labelled the model's \emph{kinematic reference}, not linear RDT
(\S\,4.1, figure~7). The linear reference against which the line model must
be judged is now the exact projected RDT, and it is computed, not asserted:
\S\,4.2 and figure~8 give the closed-form Cauchy solution for an isotropic
initial spectrum, integrated over the transverse wavenumber plane, and show
that the exactly distorted, line-projected component spectra are
\emph{not} a rigid translation---the shape ratio to a rigid translation
departs from unity and the components split. The manuscript now states
explicitly that the reference already differs from the model's own
kinematics at the linear level. The consequences the referee flags are
correspondingly revised: the claim of ``suppressed spectral migration''
relative to a rigid law is replaced by a measured comparison against exact
RDT and DNS, and the spectral results are validated against the
strained-box DNS of \S\,4.6 (figure~11, table~1), with the
distance-from-linear-theory envelope of \S\,4.6.2 (figure~12) as the
quantitative statement. The referee is also right that the \S\,4.1
(formerly Reynolds-stress-level) tests do not establish a spectral result;
that is now the explicit reason the strained-box DNS was performed
(\S\,4.6, opening paragraph).

\point{2. The claimed single-line closure of the rapid pressure--strain
term is not closed.}
\begin{refcomment}
``Eq.\ (5.7) still requires the two-variable spectral functions
$a(k_2,k_\perp)$ and $c(k_2,k_\perp)$ and integration over the unresolved
transverse wavenumber $k_\perp$. \ldots\ the paper appears to derive only
an axisymmetric 2D representation, not a single-line closure. Moreover,
this spectral formulation is not used in the simulations \ldots''
\end{refcomment}

\response
The referee's diagnosis is correct and is now the stated position of the
section, not a claim it contends with. Section~5.3 carries out the
axisymmetric collapse and then states its status precisely: under the
axisymmetry hypothesis the rapid term is a linear functional of the two
scalars $a(\kappa_2,\kappa_\perp)$ and $c(\kappa_2,\kappa_\perp)$ on the
resolved half-plane (equation~(5.6) with kernels~(5.7)--(5.9))---``a
two-dimensional representation, not yet a single-line closure, since the
line determines only the $\kappa_\perp$-integrals of fixed combinations of
$a$ and $c$ and not the functions themselves.'' The infinite-dimensional
null space of the map from line data to the rapid term is stated exactly
(\S\,5.2). What the line \emph{does} determine is then made quantitative:
because the rapid term is linear in $(a,c)$, its extreme values over
everything consistent with the line spectra, solenoidality and
realisability are the solution of a family of linear programs, giving sharp
bounds $[\Pi^-_{nn},\Pi^+_{nn}]$ and a tabulated kinematic indeterminacy
(\S\,5.4, table~2; forward map and linear programs in Appendices~A.2 and
A.4). The representation also yields two falsifiers the line reports for
itself ($\phi_{11}=\phi_{33}$ and vanishing cross-spectra), so the
hypothesis is monitored rather than assumed. On the referee's second
sentence: we agree, and now say plainly that the spectral formulation is
not exercised in the simulations---the reported results use the moment
closure to which equation~(5.6) reduces in the isotropic limit
(\S\,5.7)---and we identify the wavenumber-dependent kernel as the missing
ingredient whose implementation is future work (\S\,6.4). The section's
claim is thus reduced to what is demonstrated: an exact obstruction, sharp
bounds from line data, a measured hypothesis, and exactness for
axisymmetric distortion.

\point{3. The theoretical presentation is not sufficiently clear or
self-contained.}
\begin{refcomment}
``The manuscript frequently mixes established RDT results, existing
closures, assumptions, and new developments without clearly separating
them. In Section 5, `A1' is repeatedly used without formal definition,
while Eqs.\ (5.8)-(5.10), described as a central result, appear essentially
without derivation.''
\end{refcomment}

\response
The presentation has been reorganised to separate the four categories.
Established RDT results (the exact rapid pressure--strain integral and its
spectral form) are attributed to the classical sources and Sagaut \& Cambon
(2018) and presented as standard (\S\S\,2.4.1, 5.1); the existing algebraic
closures (Rotta, IP, LRR) are identified as such and confined to
\S\,2.4.2; the assumptions are stated as labelled hypotheses; and the new
developments (the obstruction, the bounds, the measured hypothesis, the
kernel reduction) are collected in Section~5. The hypothesis the referee
refers to, A1, is now given a formal definition at first use---``Assumption
A1 (axisymmetry of the undistorted spectrum about the line)'' in
\S\,5.3---with its justification for the present configuration and a
subsection (\S\,5.5, table~3, figure~13) that measures the error it incurs
under plane strain rather than presuming it. The kernel equations the
referee cites (now equations~(5.7)--(5.9)) are derived in Appendix~A.3 from
the azimuthal average of the exact projection weights: $g_2$ is worked
through in full and $g_1,g_3$ follow by the same steps, and the two exact
constraints that verify the reduction---the pointwise trace condition
$g_1+g_2+g_3=0$ (equation~(5.10)) and the isotropic ratio
$+\tfrac15:-\tfrac15:0$ (equation~(5.11))---are stated and checked in the
main text. We would welcome the additional issues the referee chose not to
enumerate and will address them in a further revision if the paper is
reconsidered.

\end{document}
